\section{Continuity and Convergence}



\subsection{Convergence in Topological Spaces}

We will try to answer the following questions:

\begin{enumerate}
    \item Can we define sequential convergence in general topological spaces?
    \item If so, does this convergence describe the notion of topology?
\end{enumerate}

In metric space setting, we have the notion of convergence of a sequence (in the language of open balls):

\begin{proposition}
    A sequence $(x_n)_{n\in\mathbb{N}}$ in a metric space $X$ converges to a point $x\in X$ if and only if for every $\epsilon>0$, there exists $N\in\mathbb{N}$ such that $x_n\in B_\epsilon(x)$ for all $n\ge N$.
\end{proposition}

This motivates the following definition in general topological spaces:



\begin{definition}
    A sequence $(x_n)_{n\in\mathbb{N}}$ in a topological space $X$ is said to converge to a point $x\in X$ if for every neighbourhood $U$ of $x$, there exists $N\in\mathbb{N}$ such that $x_n\in U$ for all $n\ge N$.
\end{definition}

\begin{problem}
    A sequence $(x_n)$ in a metric space converges in the metric sense if and only if it converges in the topological sense defined above.
\end{problem}
\begin{proof}
    ($\Rightarrow$) Let $(x_n)$ be a sequence in a metric space $(X,d)$ converging to $x\in X$ in the metric sense. For every neighborhood $U$ of $x$, there exists $\epsilon>0$ such that $B_\epsilon(x)\subseteq U$. Then, there exists $N$ such that $d(x_n,x)<\epsilon$ for all $n\ge N$. In particular, $x_n\in B_\epsilon(x)\subset U$ for all $n\ge N$. So we have $x_n\to x$ in the topological sense.

    ($\Leftarrow$) Let $(x_n)$ be a sequence in a metric space $(X,d)$ converging to $x\in X$ in the topological sense. For every $\epsilon>0$, the open ball $B_\epsilon(x)$ is a neighbourhood of $x$, so there exists $N$ such that $x_n\in B_\epsilon(x)$ for all $n\ge N$. This means that $d(x_n,x)<\epsilon$ for all $n\ge N$. Hence, $(x_n)$ converges to $x$ in the metric sense.
\end{proof}



\begin{proposition}
   Let $X$ be a topological space. For $A\subseteq X$ and $x\in X$, if there is a sequence $(a_n)\subseteq A$ with $a_n\to x$, then $x\in\overline{A}$.
\end{proposition}
\begin{proof}
    Let $U$ be any neighbourhood of $x$. Then there exists $N$ such that $a_n\in U$ for all $n\ge N$. In particular, $a_N\in A\cap U$, so $A\cap U\neq \emptyset$. Since $U$ was arbitrary, we have $x\in\overline{A}$.
\end{proof}

We know the converse is true in metric spaces. Is it true in general topological spaces?


% \begin{definition}
%     A topological space $X$ is first countable if every point $x\in X$ has a countable neighbourhood basis.
% \end{definition}

% Under this assumption, we have the following characterization of convergence:
% \begin{theorem}
%     If $X$ is first countable and $E\subset X$, then $x\in\overline{E}$ iff there exists a sequence $(x_n)_{n\in\mathbb{N}}$ in $E$ such that $x_n\to x$.
% \end{theorem}

% \begin{corollary}
%     Let $X$ and $Y$ be first countable topological spaces. Then,
%     \begin{enumerate}
%         \item $U\subset X$ is open iff whenever $x_n\to x\in U$, $x_n$ is eventually in $U$.
%         \item $F\subset X$ is closed iff whenever $(x_n)\subset F$, $x_n\to x$, then $x\in F$.
%         \item $f:X\to Y$ is continuous iff whenever $x_n\to x\in X$, then $f(x_n)\to f(x)\in Y$.
%     \end{enumerate}
% \end{corollary}

% For first countable spaces, the notion of sequential convergence describes the topology.


\begin{example}
Let $X=\R$ be endowed with the 
\emph{countable complement topology}: a set $U\subseteq X$ is open 
iff $X\setminus U$ is countable or $U=\emptyset$. 

\begin{itemize}
  \item Any sequence $(x_n)$ converging to $x$ is eventually equal to $x$.  
  Indeed, we let $B=\{x_n:\,x_n\neq x\}$ and hence $U=X\setminus B$ is an 
  open neighborhood of $x$. Then, there exists $N$ such that $x_n\in U$ for all
  $n\ge N$, i.e., $x_n\notin B$ for all $n\ge N$. This means that 
  $x_n=x$ for all $n\ge N$.
  \item The closed sets are exactly the countable sets and $X$ itself.  
  Hence, if $A\subsetneq X$ is uncountable, then 
  $\overline{A}=X$.
\end{itemize}

Thus $\overline{[0,1]}=\R$, but there is no sequence in $[0,1]$ converging to $2$.
\end{example}

Let's see another example.

\begin{example}
    Let $X=\R^\R$ be endowed with the product topology. We consider 
    $$E=\{f\in \R^\R:f(x)=0\text{ or }1\text{ and }f(x)=0\text{ only finitely often}\}.$$

    Let $\mathbf{0}$ denote the all 0 function in $X$. Note that any basis element of the product topology containing $\mathbf{0}$ is of the form
    $$V=\{g\in \R^\R:g(x_i)\in U_i\text{ for }i=1,\ldots,n\}$$
    for some $x_1,\ldots,x_n\in\R$ and open sets $U_i\ni\R$ containing $0$. We now define $g:\R\to \R$ as follows: $g(x_i)=0$ for $i=1,\ldots,n$ and $g(x)=0$ otherwise. Then, $g\in E\cap V$ and hence $E\cap V\neq\emptyset$. Hence, $\mathbf{0}\in\overline{E}$.
    
    Now let $(f_n)$ be a sequence in $E$ such that each $f_n$ is 0 on a finite set $A_n\subset \R$. If $(f_n)$ converges to $f$, then $f$ can be zero at most on the countable set $\bigcup_{n} A_n$. Suppose otherwise $f(x)=0$ for some $x\notin \bigcup_n A_n$. Pick the basic open neighbourhood
    $$U=\{g\in \R^\R:g(x)\in (-1,1)\}.$$
    Then, there exists $N$ such that $f_n\in U$ for all $n\ge N$. In particular, $f_n(x)\in (-1,1)$ for all $n\ge N$. This implies that $f_n(x)=0$ for all $n\ge N$, contradicting the choice of $x$.
    
    Hence, there is no sequence in $E$ converging to $\mathbf{0}$.
\end{example}

In metric spaces, the limit of a convergent sequence is unique. Is this true in general topological spaces?

\begin{example}
    Let $X$ be endowed with the trivial topology $\tau=\{\emptyset,X\}$. Then the constant sequence $(x_n)$ with $x_n=a\in X$ for all $n\in\mathbb{N}$ converges to every point in $X$. 
\end{example}

\begin{example}
    Let $X=\{a,b\}$ with the topology $\tau=\{\emptyset,\{a\},X\}$. Then the constant sequence $(x_n)$ with $x_n=a$ for all $n\in\mathbb{N}$ converges to $a$ and $b$. 
\end{example}

In general, the limit of a convergent sequence may not be unique. We will need to give some additional conditions on the topology to ensure uniqueness of limits.

Finally, let's talk about continuity.

\begin{proposition}
    Let $f:X\to Y$ be continuous at $x$. Then, whenever $x_n\to x$ in $X$, we have $f(x_n)\to f(x)$ in $Y$.
\end{proposition}
\begin{proof}
    Let $U$ be any neighbourhood of $f(x)$ in $Y$. By continuity of $f$ at $x$, there exists a neighbourhood $V$ of $x$ such that $f(V)\subseteq U$. Since $x_n \to x$, there exists $N$ such that for all $n \geq N$, $x_n \in V$. Hence, for all $n \geq N$, $f(x_n) \in U$, which shows that $f(x_n) \to f(x)$.
\end{proof}

The converse is not true in general. We will see an example below. But before that, let's see an interesting consequence.

\begin{theorem}
    Let $X=\Pi X_\alpha$ be a product of topological spaces $X_\alpha$. A sequence $(x_n)$ in $X$ converges to $x\in X$ if and only if for every $\alpha$, the sequence $(\pi_\alpha(x_n))$ converges to $\pi_\alpha(x)$ in $X_\alpha$.
\end{theorem}
\begin{proof}
    ($\Rightarrow$) This follows from the continuity of the projection maps $\pi_\alpha$.

    ($\Leftarrow$) Let $(x_n)$ be a sequence in $X$ such that for every $\alpha$, the sequence $(\pi_\alpha(x_n))$ converges to $\pi_\alpha(x)$ in $X_\alpha$. Let $U$ be a subbasis neighbourhood of $x$ in $X$ of the form  
    $$U=\pi_{\alpha}^{-1}(V)=V\times \prod_{\beta \neq \alpha} X_\beta,$$
    where $\alpha$ is an index and $V$ is a neighbourhood of $\pi_{\alpha}(x)$ in $X_{\alpha}$. 
    Since $\pi_{\alpha}(x_n)\to \pi_{\alpha}(x)$, there exists $N$ such that $\pi_{\alpha}(x_n)\in V$ for all $n\ge N$. Then, for all $n\ge N$, we have $x_n\in U$. Hence, $(x_n)$ converges to $x$ in $X$.
\end{proof}

\begin{example}[Sequential continuity does not imply continuity]
Let
\[
X=\{0,1\}^{\R}
\]
have the product topology, where each factor $\{0,1\}$ is discrete. For $x\in X$, write
\[
\operatorname{supp}(x)=\{t\in\R:x(t)=1\}.
\]
Let $S=\{0,1\}$ have the topology $\{\varnothing,\{1\},S\}$, and define
\[
f:X\longrightarrow S,\qquad
f(x)=
\begin{cases}
0,&\operatorname{supp}(x)\text{ is countable},\\
1,&\operatorname{supp}(x)\text{ is uncountable}.
\end{cases}
\]
Then $f$ preserves limits of sequences but is not continuous.

\begin{proof}
We first show that $f$ is not continuous. Let $\mathbf{1}\in X$ be the constant function with value $1$. Since $f(\mathbf{1})=1$ and $\{1\}$ is open in $S$, continuity at $\mathbf{1}$ would require a neighborhood of $\mathbf{1}$ on which $f$ is constantly $1$.

However, every basic neighborhood of $\mathbf{1}$ restricts only finitely many coordinates, say those in $F\subseteq\R$. The function
\[
y(t)=
\begin{cases}
1,&t\in F,\\
0,&t\notin F
\end{cases}
\]
belongs to that neighborhood but has finite support, so $f(y)=0$. Thus $f$ is not continuous at $\mathbf{1}$.

To prove sequential continuity, we use the following observation:
\begin{quote}
If $z_n\to z$ in $X$ and each $z_n$ has countable support, then $z$ also has countable support.
\end{quote}
Indeed, the set
\[
C=\bigcup_{n=1}^{\infty}\operatorname{supp}(z_n)
\]
is countable. For every $t\notin C$, we have $z_n(t)=0$ for all $n$. Since convergence in the product topology implies coordinatewise convergence, $z(t)=0$. Consequently,
\[
\operatorname{supp}(z)\subseteq C,
\]
which proves the observation.

Now suppose $x_n\to x$ in $X$. Recall that in $S$, every sequence converges to $0$, whereas a sequence converges to $1$ exactly when it is eventually equal to $1$.

If $f(x)=0$, we therefore have $f(x_n)\to f(x)$ automatically.

If $f(x)=1$, we claim that $f(x_n)=1$ eventually. Otherwise, there would be a subsequence $(x_{n_k})$ such that every $x_{n_k}$ has countable support. But $x_{n_k}\to x$, so the observation would imply that $x$ has countable support, contradicting $f(x)=1$.

Hence $f(x_n)\to f(x)$ in both cases.
\end{proof}
\end{example}


\subsection{Nets}
How do we fix this? We need a more general notion of convergence.

Note that a sequence $(x_n)$ in $X$ is in fact a function from $x:\mathbb{N}\to X$. Instead of just using $\mathbb{N}$ as the index set, we can use more general index sets.

\begin{definition}
    A set $\Lambda$ is a \emph{directed set} if it is a non-empty set with a relation $\le$ such that
    \begin{itemize}
        \item reflexive: for all $\lambda\in\Lambda$, $\lambda\le \lambda$,
        \item transitive: for all $\lambda_1,\lambda_2,\lambda_3\in\Lambda$, if $\lambda_1\le \lambda_2$ and $\lambda_2\le \lambda_3$, then $\lambda_1\le \lambda_3$,
        \item directed: for all $\lambda_1,\lambda_2\in\Lambda$, there exists $\mu\in\Lambda$ such that $\lambda_1\le \mu$ and $\lambda_2\le \mu$.
    \end{itemize}
    
    
    This relation $\le$ is called a direction on $\Lambda$.
\end{definition}

\begin{remark}
    Note a direction may not be a partial order, i.e., it may not be antisymmetric. For example, let $\Lambda=\{a,b\}$ with $a\le b$ and $b\le a$. This is a directed set but not a poset because $a\neq b$.
\end{remark}

\begin{example}
    $\N$ with the standard order (which is obviously a direction) is a directed set.
\end{example}


\begin{definition}
    A net in a topological space $X$ is a function $P\colon \Lambda\to X$ where $\Lambda$ is a directed set. We often write $P(\lambda)=x_\lambda$ and write the net $(x_\lambda)_{\lambda\in\Lambda}$ or $(x_\lambda)$.

    % A subnet of a net $(x_\lambda)_{\lambda\in\Lambda}$ is the composition $P\circ \varphi$ where $\varphi:M\to \Lambda$ is a map between directed sets such that 
    % \begin{enumerate}
    %      \item for every $\mu_1,\mu_2\in M$ such that $\mu_1\le\mu_2$, we have $\varphi(\mu_1)\le\varphi(\mu_2)$.
    %     \item for every $\lambda\in\Lambda$, there exists $\mu\in M$ such that $\lambda\le\varphi(\mu)$.
    % \end{enumerate}
    % For $\mu\in M$, the point $P\circ \varphi(\mu)$ is often denoted by $x_{\lambda_\mu}$ and the subnet is denoted by $(x_{\lambda_\mu})_{\mu\in M}$. (think of the case of subsequences to understand this definition)
\end{definition}

\begin{definition}
    We say that a net $(x_\lambda)_{\lambda\in\Lambda}$ converges to a point $x\in X$ if for every neighbourhood $U$ of $x$, there exists $\lambda_0\in\Lambda$ such that $x_\lambda\in U$ for all $\lambda\ge \lambda_0$.
\end{definition}

\begin{remark}
    Note that if we take $\Lambda=\N$ with the standard order, then the notion of convergence of nets reduces to the notion of convergence of sequences.
\end{remark}

\begin{example}
    Let $\mathcal{P}$ denote the collection of all finite partitions of the closed interval $[a,b]$ into closed subintervals. In other words, each $P\in \mathcal{P}$ has the form $P=\{[x_0,x_1],[x_1,x_2],\ldots,[x_{n-1},x_n]\}$ where $a=x_0<x_1<\cdots<x_n=b$.
    
    This becomes a directed set if we define the direction by refinement: for two partitions $P_1,P_2\in\mathcal{P}$, we say that $P_1\le P_2$ if $P_2$ is a refinement of $P_1$ (i.e., every subinterval in $P_1$ is a union of subintervals in $P_2$).

    Now, consider a function $f:[a,b]\to \R$. Define a net $(R_L(P))_{P\in \mathcal{P}}$ as follows: for each partition $P=\{[x_0,x_1],[x_1,x_2],\ldots,[x_{n-1},x_n]\}$, let
    $$R_L(P)=\sum_{i=1}^n \inf_{x\in [x_{i-1},x_i]} f(x) \cdot (x_i - x_{i-1}).$$
    This net is called the \emph{lower Riemann sum} net of $f$.

    Similarly, we can define the \emph{upper Riemann sum} net $(R_U(P))_{P\in \mathcal{P}}$ by
    $$R_U(P)=\sum_{i=1}^n \sup_{x\in [x_{i-1},x_i]} f(x) \cdot (x_i - x_{i-1}).$$

    If both nets converge to the same limit $L$, then we say that $f$ is Riemann integrable on $[a,b]$ with integral $L$.
\end{example}


\begin{example}
    One typical example of a net is the following: let $X$ be a topological space and $x\in X$. Let $\Lambda$ be the set of all neighbourhoods of $x$ and define the direction on $\Lambda$ by $U\le V$ iff $V\subseteq U$ (reverse inclusion). You should check that this is indeed a directed set.  Pick a point $x_U\in U$ for each neighbourhood $U$ of $x$. Then the net $(x_U)_{U\in\Lambda}$ converges to $x$. To see this, let $V$ be any neighbourhood of $x$, then for all $U\geq V$, we have $x_U\in U\subseteq V$. This shows that $x_U\to x$.
\end{example}

\begin{theorem}
    If $E\subset X$, then $x\in\overline{E}$ iff there is a net $(x_\lambda)_{\lambda\in\Lambda}$ in $E$ such that $x_\lambda\to x$.
\end{theorem}
\begin{proof}
    If $x\in \overline{E}$, then every neighbourhood $U$ of $x$ intersects $E$. Choose $x_U\in E\cap U$ for each neighbourhood $U$ of $x$. Then the net $(x_U)_{U\in\Lambda}$ is a net converging to $x$.

    Conversely, if there is a net $(x_\lambda)_{\lambda\in\Lambda}$ in $E$ such that $x_\lambda\to x$, then for every neighbourhood $U$ of $x$, there exists $\lambda_0$ such that $x_\lambda\in U$ for all $\lambda\ge \lambda_0$. In particular, $x_{\lambda_0}\in E\cap U$, so $E\cap U\neq \emptyset$. Since $U$ was arbitrary, we have $x\in \overline{E}$.
\end{proof}

\begin{theorem}
    Let $f:X\to Y$. Then $f$ is continuous at $x\in X$ iff for every net $(x_\lambda)_{\lambda\in\Lambda}$ in $X$ such that $x_\lambda\to x$, we have $f(x_\lambda)\to f(x)$.
\end{theorem}
\begin{proof}
    ($\Rightarrow$) Given a neighborhood $V$ of $f(x)$ in $Y$, by continuity of $f$ at $x$, there exists a neighbourhood $U$ of $x$ such that $f(U)\subseteq V$. Since $x_\lambda\to x$, there exists $\lambda_0$ such that $x_\lambda\in U$ for all $\lambda\ge \lambda_0$. Hence, for all $\lambda\ge \lambda_0$, we have $f(x_\lambda)\in V$, which shows that $f(x_\lambda)\to f(x)$.

    ($\Leftarrow$) Suppose $f$ is not continuous at $x$. Then there exists a neighbourhood $V$ of $f(x)$ such that for every neighbourhood $U$ of $x$, there exists $x_U\in U$ with $f(x_U)\notin V$. Let $\Lambda$ be the set of all neighbourhoods of $x$ directed by reverse inclusion. Then the net $(x_U)_{U\in\Lambda}$ converges to $x$, but $f(x_U)\notin V$ for all $U\in\Lambda$. Hence, $f(x_U)$ does not converge to $f(x)$.
\end{proof}



% TODO: include examples of nets, subnets, and their use in proofs of
% continuity and closure properties.

%---------------------------------------------------------------------------
% Section 5: Separation and countability axioms (Willard §§13–16)
%---------------------------------------------------------------------------

\newpage
